\documentclass[aspectratio=169,11pt]{beamer}
\usepackage{../../../shared/theme/esmad_beamer_theme}
\usepackage{amsmath,amssymb}
\usepackage{booktabs}
\usepackage{array}

\title[Difference-in-Differences]{Difference-in-Differences and Event Studies}
\subtitle{Identification from Differential Trends}
\date{\today}

\newcommand{\E}{\mathbb{E}}

\begin{document}

\begin{frame}
\titlepage
\end{frame}

\begin{frame}{Learning Goals}
\begin{itemize}
\item derive the canonical 2x2 Difference-in-Differences estimand;
\item state parallel trends using potential outcomes;
\item connect the design-based estimand to two-way fixed-effects regression;
\item interpret event-study coefficients before and after treatment;
\item recognize anticipation, staggered adoption, and heterogeneous treatment-effect problems;
\item use clustering, placebo tests, and pre-treatment diagnostics without overclaiming what they prove.
\end{itemize}
\end{frame}

\begin{frame}{Outline}
\tableofcontents
\end{frame}

\section{The Basic Design}

\begin{frame}{The 2x2 Difference-in-Differences Setup}
There are two groups and two periods:
\begin{itemize}
\item treated group \(G=1\);
\item comparison group \(G=0\);
\item pre-treatment period \(t=0\);
\item post-treatment period \(t=1\).
\end{itemize}

The canonical estimator is
\[
\widehat{\tau}_{DID}
=
(\bar Y_{1,1}-\bar Y_{1,0})
-
(\bar Y_{0,1}-\bar Y_{0,0}).
\]

\begin{block}{Interpretation}
Subtract the change in the untreated comparison group from the change in the treated group.
\end{block}
\end{frame}

\begin{frame}{Why the Difference of Differences?}
A before--after comparison in the treated group,
\[
\bar Y_{1,1}-\bar Y_{1,0},
\]
mixes treatment with all other time changes.

A treated--control comparison after treatment,
\[
\bar Y_{1,1}-\bar Y_{0,1},
\]
mixes treatment with persistent group differences.

\begin{alertblock}{DiD logic}
Use the comparison group's change to estimate the treated group's counterfactual change in the absence of treatment.
\end{alertblock}
\end{frame}

\section{Potential Outcomes and Identification}

\begin{frame}{Potential Outcomes with Groups and Time}
Let \(Y_{it}(1)\) denote the outcome under treatment and \(Y_{it}(0)\) the untreated potential outcome.

For treated units after intervention,
\[
Y_{i1}=Y_{i1}(1).
\]

The missing counterfactual is
\[
Y_{i1}(0).
\]

Difference-in-Differences identifies this missing counterfactual through assumptions about untreated trends.
\end{frame}

\begin{frame}{Parallel Trends}
A canonical parallel-trends condition is
\[
\E[Y_1(0)-Y_0(0)\mid G=1]
=
\E[Y_1(0)-Y_0(0)\mid G=0].
\]

\begin{block}{Meaning}
Without treatment, the average outcome change would have been the same in treated and comparison groups.
\end{block}

\begin{alertblock}{What is not required}
The groups do not need equal outcome levels before treatment.
\end{alertblock}
\end{frame}

\begin{frame}{What Parallel Trends Does Not Say}
Parallel trends does not require:
\begin{itemize}
\item identical levels;
\item identical covariate distributions;
\item identical shocks after treatment;
\item constant treatment effects;
\item linear outcome paths.
\end{itemize}

It does require the comparison group's untreated evolution to be informative about the treated group's untreated evolution over the relevant horizon.
\end{frame}

\begin{frame}{Conditional Parallel Trends}
Sometimes parallel trends is more credible after conditioning on covariates \(X\):
\[
\E[Y_1(0)-Y_0(0)\mid G=1,X]
=
\E[Y_1(0)-Y_0(0)\mid G=0,X].
\]

\begin{itemize}
\item covariates should be pre-treatment or otherwise justified;
\item overlap between treated and comparison units remains important;
\item conditioning cannot fix an omitted time-varying confounder absent supporting assumptions.
\end{itemize}
\end{frame}

\section{Regression Formulation}

\begin{frame}{The Canonical Regression}
For the 2x2 design:
\[
Y_{it}
=
\alpha
+
\gamma G_i
+
\lambda Post_t
+
\beta(G_i\times Post_t)
+
u_{it}.
\]

The interaction coefficient \(\beta\) equals the sample DiD contrast.

\begin{alertblock}{Important}
The regression representation does not eliminate the need for parallel trends.
\end{alertblock}
\end{frame}

\begin{frame}{Two-Way Fixed Effects}
With many units and periods:
\[
Y_{it}
=
\alpha_i
+
\lambda_t
+
\beta D_{it}
+
u_{it}.
\]

Unit fixed effects absorb stable unit differences. Time fixed effects absorb shocks common to all units.

\begin{block}{Identification}
The coefficient is learned from treatment-status changes after removing unit and time averages.
\end{block}
\end{frame}

\section{Event Studies}

\begin{frame}{Relative Event Time}
Let \(T_i\) be the first treatment period and define event time
\[
k=t-T_i.
\]

An event-study specification is
\[
Y_{it}
=
\alpha_i+\lambda_t
+
\sum_{k\neq -1}
\beta_k
\mathbb{1}\{t-T_i=k\}
+
u_{it}.
\]

Usually one pre-treatment period, often \(k=-1\), is omitted as the reference.
\end{frame}

\begin{frame}{How to Read Event-Study Coefficients}
\begin{itemize}
\item \(\beta_{-3},\beta_{-2}\): pre-treatment relative differences;
\item \(\beta_0\): effect around treatment onset;
\item \(\beta_1,\beta_2,\ldots\): dynamic post-treatment effects;
\item all coefficients are relative to the omitted event-time period.
\end{itemize}

\begin{alertblock}{Pre-treatment coefficients}
They can reveal obvious violations of parallel trends, but insignificant coefficients do not prove parallel trends.
\end{alertblock}
\end{frame}

\begin{frame}{Why Pre-Trend Tests Have Limited Power}
Suppose the untreated trends differ modestly but the pre-period is short or noisy.

Then:
\begin{itemize}
\item pre-treatment estimates can be imprecise;
\item a joint test may fail to reject;
\item the post-treatment bias may still be meaningful.
\end{itemize}

\begin{block}{Good practice}
Inspect magnitudes and patterns, not only p-values.
\end{block}
\end{frame}

\section{Anticipation and Dynamics}

\begin{frame}{Anticipation Violates the Simple Timing Story}
If units change behavior before formal treatment starts,
\[
Y_{it}(1)\neq Y_{it}(0)
\]
for nominally pre-treatment periods near adoption.

Examples:
\begin{itemize}
\item firms invest before a regulation takes effect;
\item customers stockpile before a price increase;
\item workers change behavior after a policy announcement.
\end{itemize}

\begin{alertblock}{Consequence}
The assumed untreated pre-period may already be contaminated by treatment.
\end{alertblock}
\end{frame}

\begin{frame}{Dynamic Treatment Effects}
Treatment effects may evolve:
\[
\tau_k
=
\E[Y_{i,T_i+k}(1)-Y_{i,T_i+k}(0)\mid G_i=1].
\]

Possible patterns include:
\begin{itemize}
\item immediate jump;
\item delayed onset;
\item cumulative growth;
\item temporary effect;
\item fade-out.
\end{itemize}

Event studies are useful precisely because a single static post-treatment coefficient can hide this structure.
\end{frame}

\section{Staggered Adoption}

\begin{frame}{Treatment at Different Times}
In staggered designs, groups adopt treatment in different periods:
\[
T_i\in\{2,3,4,\ldots\}.
\]

At a given time, already-treated units can inadvertently serve as controls for newly treated units in a standard TWFE regression.

\begin{alertblock}{Why this matters}
With heterogeneous treatment effects, those comparisons can receive unintuitive or even negative weights.
\end{alertblock}
\end{frame}

\begin{frame}{Treatment-Effect Heterogeneity}
Suppose cohort \(g\) has an effect
\[
ATT(g,t)
=
\E[Y_t(1)-Y_t(0)\mid T_i=g].
\]

If effects differ across cohorts or event time, a single TWFE coefficient need not equal a simple average of economically meaningful treatment effects.

\begin{block}{Modern strategy}
Estimate cohort- and time-specific effects using appropriate comparison groups, then aggregate transparently.
\end{block}
\end{frame}

\begin{frame}{Who Should Be the Comparison Group?}
Candidate comparisons include:
\begin{itemize}
\item never-treated units;
\item not-yet-treated units;
\item a carefully selected subset with credible untreated trends.
\end{itemize}

The comparison group should approximate the treated cohort's untreated counterfactual at that time.

\begin{alertblock}{Design principle}
Comparison-group choice is an identification decision, not a default regression setting.
\end{alertblock}
\end{frame}

\section{Inference}

\begin{frame}{Serial Correlation and Standard Errors}
DiD residuals are often correlated within treatment units over time:
\[
\Cov(u_{it},u_{is})\neq 0.
\]

Ignoring this dependence can severely understate uncertainty.

\begin{block}{Typical practice}
Cluster at the treatment-assignment level when there are sufficiently many independent clusters.
\end{block}
\end{frame}

\begin{frame}{Few Clusters}
When the number of clusters is small:
\begin{itemize}
\item conventional cluster-robust asymptotics may be poor;
\item degrees-of-freedom corrections matter;
\item randomization inference or wild-cluster bootstrap methods may be useful;
\item uncertainty should be reported conservatively.
\end{itemize}

\begin{alertblock}{Reminder}
More observations within a handful of treated clusters do not create many independent treatment assignments.
\end{alertblock}
\end{frame}

\section{Diagnostics and Falsification}

\begin{frame}{Placebo Treatment Dates}
Assign a false intervention date in the pre-treatment period and re-estimate the design.

A large placebo effect may indicate:
\begin{itemize}
\item differential pre-trends;
\item anticipation;
\item model misspecification;
\item seasonal or structural breaks.
\end{itemize}

\begin{block}{Interpretation}
A null placebo is reassuring only for the specific failure modes the placebo could expose.
\end{block}
\end{frame}

\begin{frame}{Placebo Outcomes}
Estimate the design on outcomes that should not respond to treatment.

A detected effect can signal:
\begin{itemize}
\item omitted concurrent shocks;
\item broad changes unrelated to treatment;
\item comparison-group failure.
\end{itemize}

Negative-control outcomes are especially useful when they share confounding pathways but not the treatment mechanism.
\end{frame}

\begin{frame}{Covariate and Composition Diagnostics}
Check whether treatment coincides with changes in:
\begin{itemize}
\item sample composition;
\item measurement procedures;
\item migration across groups;
\item attrition;
\item pre-treatment covariates that should be stable.
\end{itemize}

\begin{alertblock}{Threat}
A changing population can make apparent outcome trends differ even when individual-level potential outcomes do not.
\end{alertblock}
\end{frame}

\section{Applied Reasoning}

\begin{frame}{Example: Regional Pricing Policy}
Suppose some regions introduce a pricing policy in 2026 while others do not.

A DiD analysis asks:
\begin{itemize}
\item Were treated regions already on different demand trends?
\item Did policy adoption respond to anticipated demand?
\item Were there region-specific economic shocks?
\item Did neighboring untreated regions experience spillovers?
\item Did effects appear immediately or with delay?
\end{itemize}
\end{frame}

\begin{frame}{A Credible DiD Workflow}
\begin{enumerate}
\item Define treated cohorts, comparison units, and treatment timing.
\item State the causal estimand.
\item Plot raw outcome paths before modeling.
\item Defend parallel trends substantively.
\item Estimate event-time effects.
\item Inspect pre-treatment magnitudes and uncertainty.
\item Cluster uncertainty at the assignment level.
\item Run placebo dates and negative-control outcomes.
\item Address staggered adoption and heterogeneous effects explicitly.
\end{enumerate}
\end{frame}

\begin{frame}{Common Failure Modes}
\begin{itemize}
\item treating parallel trends as automatically true because pre-trend p-values are large;
\item using already-treated units as controls without considering heterogeneity;
\item ignoring anticipation;
\item clustering at the observation level instead of assignment level;
\item overlooking sample-composition changes;
\item interpreting a single TWFE coefficient as the universal treatment effect;
\item hiding dynamic effects inside one post-treatment dummy.
\end{itemize}
\end{frame}

\begin{frame}{Synthesis: What Makes DiD Credible?}
\begin{itemize}
\item DiD is a counterfactual design based on comparative trends.
\item Parallel trends concerns untreated potential outcomes, not merely observed pre-period similarity.
\item Event studies reveal dynamics and can expose some design failures.
\item Pre-trend tests are diagnostics, not proofs.
\item Staggered adoption and heterogeneous effects require estimators that respect cohort timing.
\item Inference should reflect treatment-assignment dependence.
\item Placebos and negative controls strengthen credibility when they target plausible failure modes.
\end{itemize}
\end{frame}

\begin{frame}{Selected References}
\small
\begin{itemize}
\item Angrist, J. D., and Pischke, J.-S. (2009). \textit{Mostly Harmless Econometrics}. Princeton University Press.
\item Cunningham, S. (2021). \textit{Causal Inference: The Mixtape}. Yale University Press.
\item Callaway, B., and Sant'Anna, P. H. C. (2021). Difference-in-Differences with Multiple Time Periods. \textit{Journal of Econometrics}, 225(2), 200--230.
\item Sun, L., and Abraham, S. (2021). Estimating Dynamic Treatment Effects in Event Studies with Heterogeneous Treatment Effects. \textit{Journal of Econometrics}, 225(2), 175--199.
\item Goodman-Bacon, A. (2021). Difference-in-Differences with Variation in Treatment Timing. \textit{Journal of Econometrics}, 225(2), 254--277.
\item Roth, J. (2022). Pretest with Caution: Event-Study Estimates after Testing for Parallel Trends. \textit{American Economic Review: Insights}, 4(3), 305--322.
\end{itemize}
\end{frame}

\contactslide

\end{document}
